\input notesmac.tex
\nopagenumbers
\centerline{\titlefont Tech and Tools, Lab 5 \& 6}
\bigskip
\line{Dr.~WANG Jian-Sheng \hfill  5/7 and 12/14 October 1999}
\bigskip
\noindent {\sl You might need to read the lecture notes or books not 
yet covered in class.  Programming in Fortran 90 is assumed if not
stated otherwise.}
\medskip
%--------------------------------------------------------------------
\problem {\sc Control statements and I/O} (Metcalf \& Reid, chapter~4,
exercise~2, page~73)
\item{} The first two terms of the Fibonacci series are both 1,
$F_0 = F_1 = 1$, and all subsequent terms are defined as the
sum of the preceding two terms, $F_n = F_{n-1} + F_{n-2}$. 
Write a Fortran 77 program which reads an integer value {\tt LIMIT}
and which computes and prints the first {\tt LIMIT} 
terms of the series.  Make sure your program compile with
{\tt f77} and {\tt f90}.

\problem {\sc Control statements and I/O} (Metcalf \& Reid, chapter~4,
exercise~3, page~73)
\item{} The coefficients of successive orders of the binomial expansion
are shown in the normal Pascal triangle form as
\ttbg
    1
   1 1
  1 2 1
 1 3 3 1
1 4 6 4 1
  .....
\tted
\item{} where the numbers at a given order are obtained by adding 1 at
the left and right ends and summing the two numbers at the previous order.
Write a program which reads an integer value {\tt LIMIT} and prints the
first {\tt LIMIT} lines of this Pascal triangle.  Use array sections
if you can.

\medskip
%--------------------------------------------------------------------
\problem {\sc Keyword and optional arguments}
\item{} (a) Define your own version of logarithmic function {\tt LOGB(X,
BASE)} such that the argument {\tt BASE} is optional.  If the {\tt
BASE} is omitted, {\tt LOGB(X)} means natural logarithm (base $e$), 
otherwise {\tt LOGB(X, BASE)} is logarithm in the specified base. 
You may use the intrinsic function {\tt LOG(X)} or {\tt LOG10(X)}.
\item{} (b) Test your function in a main program.  Note that an
interface specification block is required.

\medskip
%--------------------------------------------------------------------
\problem {\sc Traveling problem}
\item{} Suppose you want to visit a large number of world famous
cities, starting from Singapore.  To minimize cost, 
the total distances traveled must be a
minimum.  You want to write a Fortran program to schedule your itinerary.
You are giving an input data file (posted at Tech and Tools homepage,
{\tt http://www.cz3.nus.edu.sg/\char`~wangjs/SMA\_C6/citydata.dat})
containing a list of cities to be visited.
The data file also specifies the location of the cities in the form of 
four numbers $lat_d$ $lat_m$ $long_d$ $long_m$, specifying the 
latitude and longitude of a city in degree and minute.  For example,
Singapore is located at latitude $1^\circ\,17'$ and longitude
$103^\circ\,51'$, so the four numbers are {\tt 1 17 \ 103 51}.
The question is to find a shortest path which passes through each city
exactly once, assuming direct flight (straightest path) between cities.
\item{} (a) Write a function that returns the distance in kilometer
between two cities when the locations in latitude and longitude are
given as arguments.  Let assume that the Earth is a perfect sphere with
radius $R = 6370$ kilometers. 
\item{} (b) Use the following algorithm (Greedy Algorithm)
to find the best path: choose
the first city, {\tt "Singapore" },
as the starting point; choose the second one that is
closest to the first one; choose the third one closest to the second 
one, excluding cities already visited;
and so on until all cities are traveled. To make a complete trip 
you will return to the first city.   Output the total distance
traveled and the sequence of the cities passed.
\item{} (c) Is your answer really the absolute optimal choice?  
In other words, is the total distance a global minimum?

\bye

